\FloatBarrier
\newpage
\section{独立性}
\label{sec:prob-independence}

条件概率研究信息怎样改变概率。若知道一个事件发生，另一个事件的概率保持不变，就有了独立性的直观含义。下面用不需要相除的等式来定义它。

\subsection{事件的独立性}

\begin{definition}
若事件\(A,B\)满足
\begin{equation}\label{eq:prob-event-independence}
 P(A\cap B)=P(A)P(B),
\end{equation}
则称\(A,B\){\heiti 相互独立}。
\end{definition}

当\(P(B)>0\)时，此条件等价于\(P(A\mid B)=P(A)\)。在独立抛掷的两枚公平骰子中，第一枚出现偶数与第二枚大于4独立；若两事件都只涉及第一枚，则需要按实际事件关系计算，不能仅凭叙述中有两个事件就使用乘法。

独立性还保留于补事件。由
\[
 P(A\cap B^c)=P(A)-P(A\cap B)=P(A)(1-P(B)),
\]
得到\(A,B^c\)独立；同理\(A^c,B\)及\(A^c,B^c\)也独立。

互斥描述不能同时发生，独立描述概率的乘积分解。两个正概率事件若互斥，交事件概率为零，而概率乘积为正，因此不独立。

\subsection{多个事件的独立性}

\begin{definition}
若\(A_1,\ldots,A_n\)中任意非空子组均满足
\[
 P(A_{i_1}\cap\cdots\cap A_{i_r})=
 P(A_{i_1})\cdots P(A_{i_r}),
\]
则称这些事件{\heiti 相互独立}。若仅每两个事件独立，称为两两独立。
\end{definition}

两两独立不足以保证相互独立。公平抛两次硬币，令\(A\)为第一次正面，\(B\)为第二次正面，\(C\)为两次相同。三者的概率均为\(1/2\)，各两事件交集的概率均为\(1/4\)，但
\[
 P(A\cap B\cap C)=\frac14\ne\frac18.
\]
知道\(A,B\)同时发生，已经确定了\(C\)发生；各自单独考察时却看不出这种三者关系。

相互独立也允许将其中任意事件换成补事件。例如，若\(A,B,C\)相互独立，则
\[
 P(A\cap B\cap C^c)=P(A\cap B)-P(A\cap B\cap C)
 =P(A)P(B)P(C^c).
\]
对其他补事件重复使用这种分解，就得到一般情形。

\subsection{重复试验}

一次试验只记录成功与失败，成功概率为\(p\)，称为一次{\heiti Bernoulli试验}。独立重复\(n\)次且各次的\(p\)相同，就得到独立同机制的重复试验。

指定一条含\(k\)次成功、\(n-k\)次失败的结果序列，其概率为\(p^k(1-p)^{n-k}\)。所有恰有\(k\)次成功的序列共有\(\binom nk\)条，因而
\begin{equation}\label{eq:prob-bernoulli-count}
 P(\text{恰成功}k\text{次})=\binom nkp^k(1-p)^{n-k}.
\end{equation}
至少成功一次的概率为\(1-(1-p)^n\)，直接由“全部失败”的补事件得到。

独立性是试验机制的一项模型条件。来自同一有限盒子的不放回抽取，会改变后续盒中构成；独立重复测量则需要按测量机制判断，不能由次数相同代替判断。

\begin{example}[据2022年全国乙卷理科改编]
某棋手与甲、乙、丙各比赛一盘，各盘结果相互独立，获胜概率分别为\(0<p_1<p_2<p_3<1\)。记三盘中至少出现一段连续两盘获胜的概率为\(q\)。比较三种不同的第二盘对手，确定怎样安排使\(q\)最大。
\end{example}

\begin{exercise}
\mbox{}\par\nobreak

（1）独立重复试验5次，每次成功概率为\(0.4\)。求全部成功、至少成功一次及恰成功两次的概率。

（2）两个独立元件正常工作的概率分别为\(0.9,0.8\)。串联系统要求二者都正常，并联系统要求至少一个正常。求两种系统正常的概率。

（3）已知\(P(A)=P(B)=1/2\)，\(A,B\)独立。求\(P(A\cup B)\)及\(P(A\mid A\cup B)\)。观察获得新信息后，所用模型发生了什么变化。

（4）独立公平抛三次硬币，令\(A\)表示前两次相同，\(B\)表示后两次相同，\(C\)表示第一、三次相同。判断三事件是否两两独立、相互独立。

（5）（据2021年天津卷改编）甲、乙每次各猜一个谜语，猜对概率分别为\(5/6,3/5\)。每次两人的结果互不影响，各次活动也互不影响。一人猜对、另一人猜错时，猜对者获胜，否则平局。求甲一次获胜及三次活动中至少获胜两次的概率。
\end{exercise}

\FloatBarrier
\newpage
\subsection{条件信息下的独立性*}

在\(P(C)>0\)的条件模型中，若
\[
 P(A\cap B\mid C)=P(A\mid C)P(B\mid C),
\]
称\(A,B\)在条件\(C\)下独立。它描述的是当前条件模型中的乘积分解。

例如先等可能选择机器甲或乙，再用所选机器独立制造两件产品。甲的不合格率为\(0.1\)，乙为\(0.9\)。已知机器来源，两件是否不合格相互独立；不记录来源时，每件不合格的概率均为\(0.5\)，但两件都不合格的概率为
\[
 \frac12(0.1)^2+\frac12(0.9)^2=0.41\ne0.25.
\]
第一件不合格会提供关于机器来源的信息，也就改变了第二件不合格的概率。

反过来，独立事件在获得新信息后也可能不再独立。公平抛两次硬币，两个正面事件原先独立；已知至少一个正面时，每个正面事件的条件概率为\(2/3\)，同时发生的条件概率却为\(1/3\)。判断独立性需要明确在哪个概率模型中判断。

\begin{exercise}
在上述两台机器的模型中，已知第一件不合格，求第二件也不合格的概率；再求已知第一件合格时第二件不合格的概率。比较两者与事先的\(0.5\)。
\end{exercise}
